\documentclass[11pt,a4paper]{report}

\usepackage{fontspec}
% Police principale : Arial. Arial étant propriétaire et absente sous Linux,
% on utilise « Liberation Sans », métriquement identique à Arial.
\setmainfont{Liberation Sans}
\usepackage[french]{babel}
\usepackage{graphicx}
\usepackage{float}
\usepackage{xcolor}
\usepackage{tabularx}
\usepackage{booktabs}
\usepackage{array}
\usepackage{longtable}
\usepackage{enumitem}
\usepackage{titlesec}
\usepackage{fancyhdr}
\usepackage{geometry}
\usepackage{titling}
\usepackage[hidelinks]{hyperref}

\geometry{margin=2.4cm}

% ----- Palette (alignée sur l'identité visuelle de l'application) -----
\definecolor{primary}{HTML}{1C6FB8}
\definecolor{accent}{HTML}{409EFF}
\definecolor{darknavy}{HTML}{001529}
\definecolor{success}{HTML}{2E9E5B}
\definecolor{warn}{HTML}{C9821B}
\definecolor{danger}{HTML}{C0392B}
\definecolor{lightbg}{HTML}{EEF2F7}
\definecolor{greytext}{HTML}{5A6675}

% ----- Mise en forme des titres -----
\titleformat{\chapter}[display]
  {\normalfont\huge\bfseries\color{primary}}
  {\flushright\normalsize\color{greytext}\MakeUppercase{\chaptertitlename}\ \thechapter}
  {6pt}{\Huge}
\titlespacing*{\chapter}{0pt}{10pt}{26pt}

\titleformat{\section}
  {\normalfont\Large\bfseries\color{darknavy}}{\thesection}{0.6em}{}
\titleformat{\subsection}
  {\normalfont\large\bfseries\color{primary}}{\thesubsection}{0.6em}{}
\titleformat{\subsubsection}
  {\normalfont\normalsize\bfseries\color{darknavy}}{}{0em}{}

% ----- En-têtes / pieds de page -----
\pagestyle{fancy}
\fancyhf{}
\renewcommand{\headrulewidth}{0.4pt}
\fancyhead[L]{\small\color{greytext}Guide utilisateur}
\fancyhead[R]{\small\color{greytext}\leftmark}
\fancyfoot[C]{\small\color{greytext}\thepage}

% ----- Listes plus aérées -----
\setlist[itemize]{leftmargin=1.4em,itemsep=2pt,topsep=3pt}
\setlist[enumerate]{leftmargin=1.6em,itemsep=2pt,topsep=3pt}

% ----- Emplacement réservé pour une capture d'écran -----
% L'image est volontairement générique : la LÉGENDE décrit
% précisément la capture à insérer pour le processus concerné.
\newcommand{\capture}[1]{%
  \begin{figure}[H]
    \centering
    \includegraphics[width=0.82\textwidth]{images/placeholder.png}
    \caption{#1}
  \end{figure}%
}

% ----- Capture d'écran réelle (image encadrée + légende) -----
%  \screenshot{images/fichier.png}{Légende}
\newcommand{\screenshot}[2]{%
  \begin{figure}[H]
    \centering
    \setlength{\fboxsep}{0pt}\setlength{\fboxrule}{0.4pt}%
    \fcolorbox{greytext}{white}{\includegraphics[width=0.9\textwidth]{#1}}
    \caption{#2}
  \end{figure}%
}

% ----- Encarts (rendus simples, compatibles export Word) -----
\newcommand{\note}[1]{\par\smallskip\noindent\textbf{\textcolor{primary}{Bon à savoir —}} #1\par\smallskip}
\newcommand{\astuce}[1]{\par\smallskip\noindent\textbf{\textcolor{success}{Astuce —}} #1\par\smallskip}
\newcommand{\attention}[1]{\par\smallskip\noindent\textbf{\textcolor{danger}{Attention —}} #1\par\smallskip}

% ----- Raccourcis -----
\newcommand{\menu}[1]{\textsf{\textbf{#1}}}      % chemin de menu / libellé d'écran
\newcommand{\bouton}[1]{\textsf{[\,#1\,]}}        % bouton cliquable
\newcommand{\champ}[1]{\textit{#1}}                % nom de champ
\newcommand{\code}[1]{\texttt{#1}}                 % valeur / référence technique

\setlength{\parskip}{0.5em}
\setlength{\parindent}{0pt}

\hypersetup{
  pdftitle={Guide utilisateur — Hôpital de zone de Ménontin},
  pdfauthor={Hôpital de zone de Ménontin},
}
